\documentclass[UTF8,fontset=none,12pt]{ctexart}
\usepackage[a4paper,margin=2.2cm]{geometry}
\usepackage{fontspec}
\usepackage{xeCJK}
% 字体设置：按 Windows + TeX Live 环境设置，避免 API 与中文内容显示异常
\setmainfont{TeX Gyre Termes}
\setsansfont{TeX Gyre Heros}
\setmonofont{TeX Gyre Cursor}
\setCJKmainfont{SimSun}[BoldFont=SimHei]
\setCJKsansfont{SimHei}
\setCJKmonofont{KaiTi}
\usepackage{amsmath,amssymb}
\usepackage{booktabs,longtable,tabularx,array,multirow}
\usepackage[most]{tcolorbox}
\usepackage{xcolor}
\usepackage{enumitem}
\usepackage{url}
\usepackage{hyperref}
\usepackage{fancyhdr}
\usepackage{listings}
\hypersetup{colorlinks=true,linkcolor=blue,urlcolor=blue}
\urlstyle{tt}
\newcommand{\api}[1]{{\footnotesize\nolinkurl{#1}}}
\pagestyle{fancy}
\fancyhf{}
\lhead{嵌入式复习知识点}
\rhead{STM32 HAL}
\cfoot{\thepage}
\setlength{\headheight}{15pt}
\setlist[itemize]{leftmargin=2em,itemsep=0.25em,topsep=0.25em}
\setlist[enumerate]{leftmargin=2.3em,itemsep=0.25em,topsep=0.25em}
\renewcommand{\arraystretch}{1.25}
\emergencystretch=3em
\sloppy
\newcolumntype{L}[1]{>{\raggedright\arraybackslash}p{#1}}
\newcolumntype{Y}{>{\raggedright\arraybackslash}X}

\definecolor{myblue}{HTML}{1F4E79}
\definecolor{mygreen}{HTML}{38761D}
\definecolor{myorange}{HTML}{B45F06}
\definecolor{mygray}{HTML}{F3F6F8}

\tcbset{boxrule=0.6pt,arc=2mm,left=1.5mm,right=1.5mm,top=1mm,bottom=1mm}
\newtcolorbox{pointbox}[1]{colback=blue!4!white,colframe=myblue,title=\textbf{#1},breakable}
\newtcolorbox{apibox}[1]{colback=green!4!white,colframe=mygreen,title=\textbf{#1},breakable}
\newtcolorbox{warnbox}[1]{colback=orange!6!white,colframe=myorange,title=\textbf{#1},breakable}
\newtcolorbox{summarybox}[1]{colback=mygray,colframe=black!40,title=\textbf{#1},breakable}

\lstdefinestyle{cstyle}{
  language=C,
  basicstyle=\ttfamily\small,
  keywordstyle=\bfseries\color{blue!60!black},
  commentstyle=\itshape\color{green!45!black},
  stringstyle=\color{orange!70!black},
  numbers=left,
  numberstyle=\tiny\color{black!45},
  stepnumber=1,
  numbersep=8pt,
  frame=single,
  breaklines=true,
  columns=fullflexible,
  tabsize=2,
  showstringspaces=false
}

\title{\bfseries 嵌入式系统复习资料}
\author{杨子颉}
\date{\today}

\begin{document}
\maketitle
\tableofcontents
\section{第 1 部分：嵌入式系统概述}

\subsection{核心知识点}
\begin{pointbox}{嵌入式系统是什么}
嵌入式系统是以应用为中心、以计算机技术为基础，软硬件可裁剪，适用于特定功能、可靠性、成本、体积和功耗要求的专用计算机系统。它不是通用电脑，而是嵌入到具体设备中完成控制、采集、通信、显示、运算等任务。
\end{pointbox}

\begin{pointbox}{嵌入式系统组成}
\begin{itemize}
  \item \textbf{硬件层}：微控制器/微处理器、存储器、GPIO、定时器、ADC、通信接口、电源、传感器、执行器等。
  \item \textbf{驱动层}：GPIO 驱动、串口驱动、定时器驱动、ADC 驱动、I2C/SPI 驱动等，负责直接操作硬件。
  \item \textbf{操作系统层}：裸机系统可以没有操作系统；复杂系统可使用 RTOS，实现任务调度、定时、信号量、消息队列等功能。
  \item \textbf{应用层}：完成具体产品功能，如温度采集、按键控制、数据显示、报警、通信上传等。
\end{itemize}
\end{pointbox}

\begin{pointbox}{嵌入式系统特点}
\begin{itemize}
  \item \textbf{专用性强}：面向特定任务设计，如饮水机控制、心率检测、温度采集。
  \item \textbf{实时性要求}：需要在规定时间内响应外部事件，如按键、中断、报警、通信接收。
  \item \textbf{资源受限}：Flash、RAM、CPU 主频、IO 数量有限，程序要尽量精简。
  \item \textbf{可靠性高}：常用于长期运行场景，需要看门狗、异常处理、掉电保护等机制。
  \item \textbf{低功耗}：电池供电设备需要休眠、低功耗时钟、低功耗外设等设计。
  \item \textbf{软硬件结合紧密}：程序设计必须结合原理图、引脚、电平、时钟和外设特性。
\end{itemize}
\end{pointbox}

\subsection{STM32 微处理器/微控制器内核}
\begin{pointbox}{STM32 与 Cortex-M 内核}
STM32 是 ST 公司基于 ARM Cortex-M 内核设计的 32 位微控制器系列。常见 STM32F103 使用 Cortex-M3 内核，具有中断控制器 NVIC、SysTick 滴答定时器、低功耗模式、丰富外设接口等特点。复习时要知道：\textbf{Cortex-M 是内核，STM32 是在内核外围集成外设后的芯片系列}。
\end{pointbox}

\begin{pointbox}{STM32 常见资源}
\begin{itemize}
  \item \textbf{内核资源}：CPU、寄存器组、NVIC、SysTick、调试接口。
  \item \textbf{存储资源}：Flash 用于存程序，SRAM 用于运行变量和堆栈。
  \item \textbf{通用外设}：GPIO、EXTI、中断、定时器、ADC、DMA、RTC、看门狗。
  \item \textbf{通信外设}：USART/UART、I2C、SPI、CAN、USB 等。
\end{itemize}
\end{pointbox}

\subsection{常用基础 API 与功能}

\begin{warnbox}{易错点}
\begin{itemize}
  \item 外设不工作时，先检查外设时钟是否使能。
  \item STM32 程序不是只写主函数，还要配置时钟、GPIO、外设、中断等硬件资源。
  \item HAL 库函数封装了寄存器操作，考试问“功能”时要能说清它背后完成的硬件配置。
\end{itemize}
\end{warnbox}

\section{第 2 部分：GPIO 引脚配置与基本输入输出}

\subsection{GPIO 核心知识点}
\begin{pointbox}{GPIO 是什么}
GPIO（General Purpose Input Output，通用输入输出口）是 MCU 与外部电路连接最常用的接口，可用于控制 LED、读取按键、驱动蜂鸣器、片选外设、作为复用外设引脚等。
\end{pointbox}

\begin{pointbox}{GPIO 常见工作模式}
\begin{itemize}
  \item \textbf{输入模式}：读取外部电平，可配合上拉、下拉或浮空输入。
  \item \textbf{推挽输出}：既能输出高电平，也能输出低电平，驱动能力较强，LED 控制常用。
  \item \textbf{开漏输出}：只能主动拉低，输出高电平需要外部或内部上拉，常用于 I2C 总线、线与逻辑、电平转换。
  \item \textbf{复用功能}：引脚交给外设控制，如 USART、SPI、I2C、TIM PWM。
  \item \textbf{模拟模式}：关闭数字输入输出缓冲，ADC 输入必须配置为模拟模式。
\end{itemize}
\end{pointbox}

\begin{pointbox}{LED、按键、ADC 输入分别怎么配}
\begin{itemize}
  \item \textbf{LED}：一般配置为推挽输出，注意判断高电平点亮还是低电平点亮。
  \item \textbf{按键}：一般配置为输入模式；若按下接地，则常用上拉输入；若按下接电源，则常用下拉输入。
  \item \textbf{ADC 输入}：必须配置为模拟输入模式，不能配置为普通数字输入。
  \item \textbf{I2C 引脚}：常配置为复用开漏输出，并需要上拉电阻。
  \item \textbf{USART TX/PWM 输出}：常配置为复用推挽输出。
\end{itemize}
\end{pointbox}

\subsection{推挽输出与开漏输出区别}
\begin{pointbox}{推挽与开漏比较}
\begin{tabularx}{\linewidth}{@{}p{3cm}p{5.5cm}Y@{}}
\toprule
\textbf{项目} & \textbf{推挽输出} & \textbf{开漏输出} \\
\midrule
输出高电平 & 可主动输出高电平 & 不能主动输出高电平，需要上拉电阻 \\
输出低电平 & 可主动输出低电平 & 可主动拉低 \\
驱动能力 & 较强，适合普通数字输出 & 高电平能力取决于上拉电阻 \\
典型应用 & LED、继电器控制、普通输出、USART TX & I2C、线与、多设备共享总线、电平转换 \\
\bottomrule
\end{tabularx}
\end{pointbox}

\subsection{GPIO 常用 API 与功能}

\begin{warnbox}{易错点}
\begin{itemize}
  \item GPIO 端口时钟没开，后面初始化和读写都可能无效。
  \item 按键输入要结合外部电路判断按下是高电平还是低电平。
  \item ADC 引脚不要配置成上拉输入，应配置为模拟模式。
\end{itemize}
\end{warnbox}

\section{第 3 部分：中断、NVIC 与 EXTI}

\subsection{中断基本概念}
\begin{pointbox}{中断是什么}
中断是 MCU 响应外部事件或内部外设事件的一种机制。当事件发生时，CPU 暂停当前程序，转去执行中断服务函数，处理完成后再返回原程序继续执行。常见中断来源包括按键外部中断、定时器溢出中断、串口接收中断、ADC 转换完成中断等。
\end{pointbox}

\begin{pointbox}{NVIC 作用}
NVIC（Nested Vectored Interrupt Controller，嵌套向量中断控制器）用于管理 Cortex-M 内核中断，包括中断使能、优先级设置、中断嵌套、挂起状态等。\textbf{优先级数值越小，优先级越高}。
\end{pointbox}

\subsection{中断优先级与分组}
\begin{pointbox}{抢占优先级与响应优先级}
\begin{itemize}
  \item \textbf{抢占优先级}：决定一个中断能否打断另一个正在执行的中断。
  \item \textbf{响应优先级/子优先级}：当多个同抢占优先级中断同时等待时，决定先响应哪一个。
  \item \textbf{优先级判断}：先比较抢占优先级，数值小者先执行；抢占优先级相同再比较子优先级，数值小者先执行。
  \item \textbf{优先级分组}：把可用优先级位划分给抢占优先级和子优先级，不同分组会影响可设置的级数。
\end{itemize}
\end{pointbox}

\subsection{EXTI 中断线映射}
\begin{pointbox}{外部中断线连接关系}
外部 GPIO 中断通过 EXTI 线进入 NVIC。一般 \api{Px0} 对应 \api{EXTI0}，\api{Px1} 对应 \api{EXTI1}，以此类推到 \api{EXTI15}。同一条 EXTI 线同一时间只能选择一个端口来源，例如 PA0 和 PB0 都对应 EXTI0，不能同时作为两个独立 EXTI0 使用。STM32F1 中常通过 AFIO 配置 EXTI 端口映射。
\end{pointbox}

\begin{pointbox}{常见 IRQ 通道}
\begin{itemize}
  \item \api{EXTI0_IRQn}：对应 EXTI0。
  \item \api{EXTI1_IRQn}：对应 EXTI1。
  \item \api{EXTI2_IRQn}：对应 EXTI2。
  \item \api{EXTI3_IRQn}：对应 EXTI3。
  \item \api{EXTI4_IRQn}：对应 EXTI4。
  \item \api{EXTI9_5_IRQn}：EXTI5 到 EXTI9 共用一个中断入口。
  \item \api{EXTI15_10_IRQn}：EXTI10 到 EXTI15 共用一个中断入口。
\end{itemize}
\end{pointbox}

\subsection{中断响应步骤与 HAL 处理流程}
\begin{pointbox}{中断响应步骤}
\begin{enumerate}
  \item 外设或外部引脚产生中断请求，中断挂起标志置位。
  \item NVIC 判断该中断是否使能，并比较优先级。
  \item CPU 保存现场，把当前寄存器压栈。
  \item CPU 根据中断向量表跳转到对应中断服务函数。
  \item 在中断服务函数中调用 HAL 处理函数并清除中断标志。
  \item 执行用户回调函数，完成具体任务。
  \item 中断返回，CPU 恢复现场，继续执行原程序。
\end{enumerate}
\end{pointbox}

\begin{pointbox}{HAL 中断处理流程}
HAL 库通常采用“中断入口函数 $\rightarrow$ HAL 统一处理函数 $\rightarrow$ 用户回调函数”的结构。例如外部中断常见流程为：
\begin{center}
\api{EXTI_IRQHandler()} $\rightarrow$ \api{HAL_GPIO_EXTI_IRQHandler()} $\rightarrow$ \api{HAL_GPIO_EXTI_Callback()}
\end{center}
这种方式的优点是代码结构统一、可移植性较好、用户只需在回调函数中写业务逻辑。
\end{pointbox}

\subsection{代码模板：按键控制 LED（查询方式）}
\begin{lstlisting}[style=cstyle]
while (1)
{
    if (HAL_GPIO_ReadPin(KEY_GPIO_Port, KEY_Pin) == GPIO_PIN_RESET)
    {
        HAL_Delay(20);   // 消抖

        if (HAL_GPIO_ReadPin(KEY_GPIO_Port, KEY_Pin) == GPIO_PIN_RESET)
        {
            HAL_GPIO_TogglePin(LED_GPIO_Port, LED_Pin);

            // 等待按键松开，防止一次按下触发多次
            while (HAL_GPIO_ReadPin(KEY_GPIO_Port, KEY_Pin) == GPIO_PIN_RESET);
        }
    }
}
\end{lstlisting}

\begin{warnbox}{易错点}
\begin{itemize}
  \item 中断优先级数字越小，优先级越高。
  \item 只配置 GPIO 中断模式还不够，还要使能 NVIC 对应 IRQ。
  \item 共用中断入口的 EXTI 线，要在中断处理函数中判断具体是哪一条线触发。
\end{itemize}
\end{warnbox}

\section{第 4 部分：定时器基础与输出比较}

\subsection{定时器核心知识点}
\begin{pointbox}{定时器作用}
定时器本质上是一个按时钟脉冲自动计数的硬件计数器。它可用于定时延时、周期中断、PWM 输出、输入捕获、输出比较、编码器接口、测频测速等功能。
\end{pointbox}

\begin{pointbox}{三种计数模式}
\begin{itemize}
  \item \textbf{向上计数}：CNT 从 0 计到 ARR，产生更新事件后回到 0。
  \item \textbf{向下计数}：CNT 从 ARR 计到 0，产生更新事件后重新装载。
  \item \textbf{中央对齐计数}：CNT 先向上计数再向下计数，常用于对称 PWM，能改善电机控制或功率控制中的波形特性。
\end{itemize}
\end{pointbox}

\subsection{时基单元构成及作用}
\begin{pointbox}{时基单元}
\begin{itemize}
  \item \textbf{PSC 预分频器}：对输入定时器时钟分频，降低计数频率。
  \item \textbf{CNT 计数器}：保存当前计数值。
  \item \textbf{ARR 自动重装载寄存器}：决定计数上限或周期。
  \item \textbf{更新事件}：CNT 溢出或到达重装载条件时产生，可触发中断、DMA 或更新寄存器。
  \item \textbf{CCR 捕获/比较寄存器}：用于输出比较、PWM 占空比或输入捕获值。
\end{itemize}
\end{pointbox}

\begin{pointbox}{定时周期计算}
若定时器输入时钟为 $f_{TIM}$，预分频值为 $PSC$，自动重装载值为 $ARR$，则定时器更新频率为：
\[
  f_{update}=\frac{f_{TIM}}{(PSC+1)(ARR+1)}
\]
定时周期为：
\[
  T=\frac{(PSC+1)(ARR+1)}{f_{TIM}}
\]
\end{pointbox}

\subsection{基本定时器与通用定时器}
\begin{pointbox}{不同定时器功能比较}
\begin{tabularx}{\linewidth}{@{}p{3cm}p{5.4cm}Y@{}}
\toprule
\textbf{类型} & \textbf{主要功能} & \textbf{典型用途} \\
\midrule
基本定时器 & 只有基本计数、更新事件，部分芯片可触发 DAC & 简单定时、周期中断、触发事件 \\
通用定时器 & 基本计数、输入捕获、输出比较、PWM、编码器接口等 & PWM、测频、测速、脉宽测量、周期任务 \\
高级定时器 & 在通用定时器基础上增加互补输出、死区、刹车等 & 电机控制、逆变器、功率驱动 \\
\bottomrule
\end{tabularx}
\end{pointbox}

\subsection{输出比较与 PWM 基础}
\begin{pointbox}{输出比较功能}
输出比较是把计数器 CNT 与比较寄存器 CCR 的值进行比较，当两者满足设定关系时，定时器改变输出电平或产生中断。PWM 就是输出比较功能的典型应用。
\end{pointbox}

\begin{pointbox}{PWM 中 ARR 与 CCR 的作用}
\begin{itemize}
  \item \textbf{ARR}：决定 PWM 周期。
  \item \textbf{PSC}：决定计数速度。
  \item \textbf{CCR}：决定有效电平持续时间，即占空比。
  \item \textbf{占空比}：$Duty=CCR/(ARR+1)\times100\%$。
\end{itemize}
\end{pointbox}

\subsection{代码模板：定时器周期中断}
\begin{lstlisting}[style=cstyle]
HAL_TIM_Base_Start_IT(&htim2);

void TIM2_IRQHandler(void)
{
    HAL_TIM_IRQHandler(&htim2);
}

void HAL_TIM_PeriodElapsedCallback(TIM_HandleTypeDef *htim)
{
    if (htim->Instance == TIM2)
    {
        HAL_GPIO_TogglePin(LED_GPIO_Port, LED_Pin);
    }
}
\end{lstlisting}

\section{第 5 部分：串行通信与 DMA}

\subsection{同步通信与异步通信}
\begin{pointbox}{同步与异步区别}
\begin{tabularx}{\linewidth}{@{}p{3cm}p{5.5cm}Y@{}}
\toprule
\textbf{项目} & \textbf{同步通信} & \textbf{异步通信} \\
\midrule
时钟线 & 通常有独立时钟线 & 通常没有独立时钟线 \\
同步方式 & 发送方提供或共享时钟 & 双方约定波特率，通过起始位和停止位同步 \\
典型接口 & SPI、I2C & UART/USART 异步模式 \\
优点 & 速度较高，时序明确 & 接线少，使用简单，适合远距离低速通信 \\
\bottomrule
\end{tabularx}
\end{pointbox}

\subsection{异步串口 UART/USART}
\begin{pointbox}{波特率与字符帧格式}
\begin{itemize}
  \item \textbf{波特率}：单位时间内传输的符号数，常见为 9600、115200 等。
  \item \textbf{字符帧}：通常由 1 位起始位、数据位、可选校验位、停止位组成。
  \item \textbf{常见格式}：115200-8-N-1，表示 115200 波特率、8 位数据位、无校验、1 位停止位。
  \item \textbf{连接方式}：A 板 TX 接 B 板 RX，A 板 RX 接 B 板 TX，双方需要共地。
\end{itemize}
\end{pointbox}

\begin{pointbox}{常见串口标志位}
\begin{itemize}
  \item \api{TXE}：发送数据寄存器空，可以写入新数据。
  \item \api{TC}：发送完成，表示移位寄存器中的数据也发送结束。
  \item \api{RXNE}：接收数据寄存器非空，表示收到数据可读。
  \item \api{IDLE}：总线空闲标志，常配合 DMA 接收不定长数据。
  \item \api{ORE}：溢出错误，表示旧数据未及时读取又来了新数据。
\end{itemize}
\end{pointbox}

\subsection{I2C 通信}
\begin{pointbox}{I2C 线数与作用}
I2C 通常使用两根线：\api{SCL} 为串行时钟线，\api{SDA} 为串行数据线。两根线一般为开漏结构，需要上拉电阻。I2C 支持一主多从，通过从机地址选择通信对象。
\end{pointbox}

\begin{pointbox}{I2C 协议要点}
\begin{itemize}
  \item \textbf{起始信号}：SCL 为高电平时，SDA 从高变低。
  \item \textbf{停止信号}：SCL 为高电平时，SDA 从低变高。
  \item \textbf{数据有效性}：SCL 高电平期间 SDA 必须保持稳定；SCL 低电平期间 SDA 才允许变化。
  \item \textbf{应答 ACK}：接收方在第 9 个时钟周期拉低 SDA 表示应答。
\end{itemize}
\end{pointbox}

\subsection{SPI 通信}
\begin{pointbox}{SPI 线数与作用}
SPI 常用四根线：\api{SCK} 为时钟线，\api{MOSI} 为主机输出从机输入，\api{MISO} 为主机输入从机输出，\api{NSS/CS} 为片选信号。SPI 可全双工通信，速度通常高于 I2C。
\end{pointbox}

\begin{pointbox}{SPI 四种时序模式}
SPI 模式由 \api{CPOL} 和 \api{CPHA} 决定：
\begin{itemize}
  \item \textbf{Mode 0}：CPOL=0，CPHA=0，空闲低电平，第一个边沿（上升沿）采样。
  \item \textbf{Mode 1}：CPOL=0，CPHA=1，空闲低电平，第二个边沿（下降沿）采样。
  \item \textbf{Mode 2}：CPOL=1，CPHA=0，空闲高电平，第一个边沿（下降沿）采样。
  \item \textbf{Mode 3}：CPOL=1，CPHA=1，空闲高电平，第二个边沿（上升沿）采样。
\end{itemize}
主从设备的 CPOL、CPHA 必须一致，否则会出现数据错位。
\end{pointbox}

\subsection{DMA 的作用}
\begin{pointbox}{DMA 是什么}
DMA（Direct Memory Access，直接存储器访问）可以在不占用 CPU 大量参与的情况下，在外设与内存、内存与内存之间搬运数据。它的核心作用是\textbf{降低 CPU 负担，提高连续数据传输效率}。ADC 连续采样、串口大量收发、SPI 屏幕刷新等场景常用 DMA。
\end{pointbox}

\begin{warnbox}{易错点}
\begin{itemize}
  \item UART 通信时 TX 要接对方 RX，RX 要接对方 TX。
  \item I2C 的 SDA、SCL 是开漏结构，通常必须有上拉电阻。
  \item SPI 主从两端的 CPOL 和 CPHA 必须一致。
  \item DMA 只是搬运数据，外设本身仍需要正确初始化和启动。
\end{itemize}
\end{warnbox}

\section{第 6 部分：A/D 转换 ADC}

\subsection{核心知识点}
\begin{pointbox}{ADC 是什么}
ADC（Analog to Digital Converter，模数转换器）把连续变化的模拟电压转换成离散的数字量。STM32 内部 ADC 常用于采集电压、电流、温度、传感器输出等信号。
\end{pointbox}

\begin{pointbox}{常考技术指标}
\begin{itemize}
  \item \textbf{分辨率}：ADC 能输出的数字位数，例如 12 位 ADC 输出范围为 $0\sim 4095$。位数越高，理论上能区分的电压越细。
  \item \textbf{量化单位/最小分辨电压}：
  \[
    \Delta V=\frac{V_{ref}}{2^N-1}
  \]
  其中 $N$ 为 ADC 位数，$V_{ref}$ 为参考电压。
  \item \textbf{转换精度}：实际转换值与真实模拟值的接近程度，受参考电压、噪声、非线性误差、采样时间等影响。
  \item \textbf{转换时间}：完成一次采样和转换所需时间，通常由采样周期和转换周期共同决定。
  \item \textbf{输入范围}：通常为 $0\sim V_{ref}$，输入电压不能超过芯片允许范围。
  \item \textbf{采样时间}：采样保持电容对输入信号采样的时间。信号源阻抗较大时，应适当加长采样时间。
\end{itemize}
\end{pointbox}

\begin{pointbox}{A/D 转换基本步骤}
\begin{enumerate}
  \item 配置 ADC 时钟、GPIO 模拟输入模式、ADC 参数和通道。
  \item 启动 ADC 转换。
  \item 等待转换完成标志置位，或使用中断/DMA 等方式等待完成。
  \item 读取 ADC 转换结果寄存器对应的数字值。
  \item 按公式把数字量换算为实际电压或物理量。
\end{enumerate}
\end{pointbox}

\subsection{ADC 工作原理与特点}
\begin{pointbox}{工作原理}
STM32 常用逐次逼近型 ADC。它会把输入模拟电压与内部 DAC 产生的比较电压逐位比较，逐步确定二进制结果。逐次逼近型 ADC 速度较快、功耗较低，适合 MCU 内部集成。
\end{pointbox}

\begin{pointbox}{常见转换方式}
\begin{itemize}
  \item \textbf{单次转换}：启动一次只采集一次，适合偶尔读取传感器。
  \item \textbf{连续转换}：一次启动后持续采集，适合实时监测。
  \item \textbf{扫描模式}：按顺序转换多个通道。
  \item \textbf{规则通道}：普通采样序列，最常用。
  \item \textbf{注入通道}：优先级更高，适合需要打断规则转换的采样任务。
  \item \textbf{查询方式}：CPU 等待转换完成，简单但占用 CPU。
  \item \textbf{中断方式}：转换完成后进入回调函数，CPU 占用较少。
  \item \textbf{DMA 方式}：转换结果自动搬运到内存，适合连续高速采样。
\end{itemize}
\end{pointbox}

\subsection{ADC 转换计算}
\begin{pointbox}{数字量换算电压}
若 ADC 分辨率为 $N$ 位，参考电压为 $V_{ref}$，读取值为 $D$，则输入电压近似为：
\[
  V_{in}=\frac{D}{2^N-1}\times V_{ref}
\]
12 位 ADC 常写为：
\[
  V_{in}=\frac{D}{4095}\times V_{ref}
\]
例如 $V_{ref}=3.3\text{ V}$，$D=2048$，则：
\[
  V_{in}\approx \frac{2048}{4095}\times 3.3\approx 1.65\text{ V}
\]
\end{pointbox}

\begin{warnbox}{易错点}
\begin{itemize}
  \item 12 位 ADC 最大值是 4095，不是 4096。公式分母通常写 $2^N-1$。
  \item ADC 引脚要配置为\textbf{模拟输入模式}，不要配置成普通输入或上拉输入。
  \item ADC 输入不能超过参考电压和芯片引脚允许电压，否则可能损坏芯片。
  \item 采样源阻抗过大时，采样时间太短会导致读数偏差。
\end{itemize}
\end{warnbox}

\section{第 7 部分：RTC 与看门狗}

\subsection{RTC 实时时钟}
\begin{pointbox}{RTC 作用}
RTC（Real-Time Clock，实时时钟）用于记录实际日期和时间，可实现年月日、时分秒、闹钟、时间戳、低功耗唤醒等功能。它常配合备份域和备用电池使用，使系统掉电后仍能保持时间。
\end{pointbox}

\begin{pointbox}{RTC 工作时钟源}
常见 RTC 时钟源有：
\begin{itemize}
  \item \textbf{LSE}：外部低速晶振，典型频率 32.768 kHz，精度较高，RTC 最常用。
  \item \textbf{LSI}：内部低速 RC 振荡器，成本低，但精度较差。
  \item \textbf{HSE 分频}：由外部高速时钟分频得到，部分型号支持，但掉电保持和低功耗场景不如 LSE 常用。
\end{itemize}
\end{pointbox}

\subsection{看门狗}
\begin{pointbox}{看门狗作用}
看门狗用于防止程序跑飞或死机。程序正常运行时需要定期“喂狗”；如果程序异常导致超时未喂狗，看门狗会复位系统，提高系统可靠性。
\end{pointbox}

\begin{pointbox}{IWDG 与 WWDG 区别}
	\begin{tabularx}{\linewidth}{@{}p{3cm}p{5.5cm}Y@{}}
\toprule
\textbf{项目} & \textbf{IWDG 独立看门狗} & \textbf{WWDG 窗口看门狗} \\
\midrule
时钟源 & LSI 内部低速时钟，独立于主时钟 & APB 时钟分频，依赖系统时钟 \\
特点 & 独立性强，适合系统级故障恢复 & 要求在规定时间窗口内喂狗，检测更严格 \\
是否能关闭 & 一旦启动通常不能停止，直到复位 & 可通过配置控制 \\
典型用途 & 防止程序死机、跑飞 & 检测过早或过晚喂狗等时序异常 \\
\bottomrule
\end{tabularx}
\end{pointbox}

\section{第 8 部分：时钟树与滴答定时器 SysTick}

\subsection{时钟树知识点}
\begin{pointbox}{系统时钟源}
STM32 时钟树用于把不同振荡源经过选择、倍频、分频后提供给 CPU、总线和外设。常见系统时钟源包括：
\begin{itemize}
  \item \textbf{HSI}：内部高速 RC 时钟，启动快、无需外部器件，但精度一般。
  \item \textbf{HSE}：外部高速晶振或外部时钟，精度较高。
  \item \textbf{PLL}：锁相环，可把 HSI 或 HSE 倍频，提高系统主频。
\end{itemize}
\end{pointbox}

\begin{pointbox}{总线与外设时钟}
\begin{itemize}
  \item \textbf{SYSCLK}：系统时钟，作为 CPU 和总线时钟的来源。
  \item \textbf{HCLK}：AHB 总线时钟，通常给内核、存储器、DMA 等使用。
  \item \textbf{PCLK1}：APB1 外设时钟，常给低速外设使用，如 TIM2、USART2、I2C 等。
  \item \textbf{PCLK2}：APB2 外设时钟，常给高速外设使用，如 GPIO、ADC、USART1、TIM1 等。
  \item \textbf{定时器时钟}：当 APB 分频系数不为 1 时，部分 STM32 定时器时钟可能为 PCLK 的 2 倍。
\end{itemize}
\end{pointbox}

\subsection{PLL 锁相环}
\begin{pointbox}{PLL 的作用}
PLL 用于把输入时钟倍频，得到更高的系统时钟。例如 STM32F103 常见配置为 HSE = 8 MHz，PLL 倍频 9，得到 SYSCLK = 72 MHz。
\[
  f_{SYSCLK}=f_{PLL\_IN}\times PLLMUL
\]
\end{pointbox}

\subsection{SysTick 滴答定时器}
\begin{pointbox}{SysTick 作用}
SysTick 是 Cortex-M 内核自带的 24 位递减计数器，常用于产生 1 ms 系统节拍。HAL 库中的 \api{HAL_Delay()} 和 \api{HAL_GetTick()} 通常依赖 SysTick。
\end{pointbox}

\begin{warnbox}{SysTick 常考点}
\begin{itemize}
  \item SysTick 属于内核外设，不属于普通通用定时器 TIMx。
  \item HAL 工程中通常每 1 ms 进入一次 \api{SysTick_Handler()}。
  \item 中断服务函数中调用 \api{HAL_IncTick()}，系统毫秒计数加 1。
  \item 使用 RTOS 时，SysTick 也可能作为操作系统节拍源。
\end{itemize}
\end{warnbox}

\section{第 9 部分：ADC 采集并通过串口打印}

\subsection{实现流程}
\begin{pointbox}{ADC 采集串口打印流程}
\begin{enumerate}
  \item 配置 ADC 引脚为模拟输入，配置 ADC 通道和采样时间。
  \item 配置 USART 的 TX/RX 引脚复用功能、波特率、数据位、停止位和校验位。
  \item 调用 \api{HAL_ADC_Start()} 启动采样。
  \item 调用 \api{HAL_ADC_PollForConversion()} 等待转换完成。
  \item 调用 \api{HAL_ADC_GetValue()} 读取 ADC 值。
  \item 用公式换算成电压或温度。
  \item 用 \api{sprintf()} 格式化字符串，再用 \api{HAL_UART_Transmit()} 发送。
\end{enumerate}
\end{pointbox}

\subsection{串口发送与接收}
\begin{pointbox}{串口配置要点}
\begin{itemize}
  \item TX 为发送脚，RX 为接收脚。两块板通信时，一般 TX 接对方 RX，RX 接对方 TX，共地。
  \item 常用配置：115200 波特率、8 位数据位、1 位停止位、无校验，即 115200-8-N-1。
  \item 查询发送适合简单打印；中断或 DMA 发送适合减少 CPU 阻塞。
  \item 若用 \api{printf()} 重定向到串口，MDK 中常需要勾选 \textbf{Use MicroLIB}，并实现 \api{fputc()}。
\end{itemize}
\end{pointbox}

\subsection{片上温度采集}
\begin{pointbox}{片上温度通道}
片上温度传感器通常使用 \api{ADC_CHANNEL_TEMPSENSOR}。以部分 STM32F1 为例，内部温度换算常见公式为：
\[
  T=\frac{V_{25}-V_{sense}}{Avg\_Slope}+25
\]
其中 $V_{25}$ 为 25 摄氏度下传感器输出电压，$Avg\_Slope$ 为平均斜率，具体参数应以对应芯片手册为准。对 STM32F103 常见典型值可写成：
\[
  T\approx \frac{1.43 - V_{sense}}{0.0043}+25
\]
\end{pointbox}

\subsection{代码模板：ADC 采集电压并串口打印}
\begin{lstlisting}[style=cstyle]
uint32_t adc_value;
float voltage;
char msg[64];

HAL_ADC_Start(&hadc1);
if (HAL_ADC_PollForConversion(&hadc1, 10) == HAL_OK)
{
    adc_value = HAL_ADC_GetValue(&hadc1);
    voltage = adc_value * 3.3f / 4095.0f;
    sprintf(msg, "ADC=%lu, V=%.2fV\r\n", adc_value, voltage);
    HAL_UART_Transmit(&huart1, (uint8_t *)msg, strlen(msg), 100);
}
HAL_ADC_Stop(&hadc1);
\end{lstlisting}

\section{第 10 部分：按键控制指示灯}

\subsection{接口电路与工作模式}
\begin{pointbox}{LED 接口电路}
LED 通常由 GPIO 口、限流电阻、发光二极管组成。常见连接方式有：
\begin{itemize}
  \item \textbf{高电平点亮}：GPIO 输出高电平，电流从 GPIO 流向 LED 和 GND。
  \item \textbf{低电平点亮}：LED 接到电源端，GPIO 输出低电平时灌电流，LED 点亮。
\end{itemize}
答题时要先判断电路是高电平有效还是低电平有效。
\end{pointbox}

\begin{pointbox}{按键接口电路}
按键常用上拉或下拉输入：
\begin{itemize}
  \item \textbf{上拉输入}：未按下为高电平，按下接地为低电平。
  \item \textbf{下拉输入}：未按下为低电平，按下接电源为高电平。
  \item \textbf{消抖}：按键机械抖动会造成多次触发，可用延时消抖、状态机消抖或定时扫描消抖。
\end{itemize}
\end{pointbox}

\subsection{GPIO 配置要点}
\begin{pointbox}{GPIO 模式选择}
\begin{itemize}
  \item LED：通常配置为\textbf{推挽输出}，输出能力强，可输出高/低电平。
  \item 按键：根据电路配置为\textbf{上拉输入}、\textbf{下拉输入}或\textbf{浮空输入}。
  \item ADC 输入：应配置为\textbf{模拟输入}，避免数字输入缓冲影响采样。
  \item 复用功能：串口、PWM 等外设控制引脚时需要配置为复用推挽或复用开漏。
\end{itemize}
\end{pointbox}

\subsection{代码模板：按键控制 LED}
\begin{lstlisting}[style=cstyle]
if (HAL_GPIO_ReadPin(KEY_GPIO_Port, KEY_Pin) == GPIO_PIN_RESET)
{
    HAL_Delay(20);
    if (HAL_GPIO_ReadPin(KEY_GPIO_Port, KEY_Pin) == GPIO_PIN_RESET)
    {
        HAL_GPIO_TogglePin(LED_GPIO_Port, LED_Pin);
        while (HAL_GPIO_ReadPin(KEY_GPIO_Port, KEY_Pin) == GPIO_PIN_RESET)
        {
            /* wait for release */
        }
    }
}
\end{lstlisting}

\section{第 11 部分：PWM 原理与控制函数}

\subsection{PWM 基本原理}
\begin{pointbox}{PWM 是什么}
PWM（Pulse Width Modulation，脉宽调制）通过改变高电平持续时间占整个周期的比例来控制等效输出。常用于 LED 调光、电机调速、蜂鸣器、舵机控制、功率变换器等。
\end{pointbox}

\begin{pointbox}{周期、频率、占空比}
定时器 PWM 的常用计算公式：
\[
  f_{PWM}=\frac{f_{TIM}}{(PSC+1)(ARR+1)}
\]
\[
  T_{PWM}=\frac{1}{f_{PWM}}
\]
\[
  Duty=\frac{CCR}{ARR+1}\times 100\% 
\]
其中 $f_{TIM}$ 为定时器输入时钟，$PSC$ 为预分频值，$ARR$ 为自动重装载值，$CCR$ 为比较值。
\end{pointbox}

\begin{pointbox}{PWM 输出比较原理}
定时器计数器 CNT 从 0 计数到 ARR。当 CNT 与 CCR 比较时，输出电平按 PWM 模式发生变化。CCR 越大，高电平持续时间越长，占空比越大。
\end{pointbox}

\subsection{PWM 参数计算例题}
\begin{pointbox}{例：72 MHz 定时器时钟，产生 1 kHz PWM}
假设 $f_{TIM}=72\text{ MHz}$，希望 $f_{PWM}=1\text{ kHz}$。可取 $PSC=71$，则定时器计数频率为：
\[
  \frac{72\text{ MHz}}{71+1}=1\text{ MHz}
\]
要得到 1 kHz：
\[
  ARR+1=\frac{1\text{ MHz}}{1\text{ kHz}}=1000
\]
所以 $ARR=999$。若占空比为 50\%，则：
\[
  CCR=(ARR+1)\times 50\%=500
\]
\end{pointbox}

\subsection{代码模板：设置 PWM 占空比}
\begin{lstlisting}[style=cstyle]
void PWM_SetDuty(TIM_HandleTypeDef *htim, uint32_t channel, float duty)
{
    uint32_t arr;
    uint32_t ccr;

    if (duty < 0.0f) duty = 0.0f;
    if (duty > 100.0f) duty = 100.0f;

    arr = __HAL_TIM_GET_AUTORELOAD(htim);
    ccr = (uint32_t)((arr + 1) * duty / 100.0f);
    __HAL_TIM_SET_COMPARE(htim, channel, ccr);
}

HAL_TIM_PWM_Start(&htim3, TIM_CHANNEL_1);
PWM_SetDuty(&htim3, TIM_CHANNEL_1, 50.0f);
\end{lstlisting}

\subsection{控制执行函数：判断超限报警}
\begin{pointbox}{常见综合题思路}
综合题可能要求“采集 ADC 值，换算物理量，判断是否超限，控制 LED/蜂鸣器/PWM”。答题时可按以下顺序写：
\begin{enumerate}
  \item 采集：\api{HAL_ADC_Start()}、\api{HAL_ADC_PollForConversion()}、\api{HAL_ADC_GetValue()}。
  \item 换算：数字量换成电压或温度、电流等物理量。
  \item 判断：与阈值比较。
  \item 执行：超限则 LED 亮、蜂鸣器响、PWM 占空比改变；未超限则恢复正常。
  \item 打印：用串口输出状态，方便调试。
\end{enumerate}
\end{pointbox}

\begin{lstlisting}[style=cstyle]
void Control_Alarm(float value, float limit)
{
    if (value >= limit)
    {
        HAL_GPIO_WritePin(LED_GPIO_Port, LED_Pin, GPIO_PIN_SET);
        PWM_SetDuty(&htim3, TIM_CHANNEL_1, 80.0f);
    }
    else
    {
        HAL_GPIO_WritePin(LED_GPIO_Port, LED_Pin, GPIO_PIN_RESET);
        PWM_SetDuty(&htim3, TIM_CHANNEL_1, 0.0f);
    }
}
\end{lstlisting}

\section{第 1--11 部分 API 总表}
\begin{center}
\small
\begin{longtable}{@{}p{2.4cm}p{5.5cm}p{7.5cm}@{}}
\toprule
\textbf{模块} & \textbf{常用 API} & \textbf{主要功能} \\
\midrule
\endfirsthead
\toprule
\textbf{模块} & \textbf{常用 API} & \textbf{主要功能} \\
\midrule
\endhead
概述 & \api{HAL_Init()} & 初始化 HAL 库与系统节拍等基础环境 \\
概述 & \api{SystemClock_Config()} & 配置系统时钟树，常由 CubeMX 生成 \\
GPIO & \api{HAL_GPIO_Init()} & 初始化引脚模式、上下拉和速度 \\
GPIO & \api{HAL_GPIO_ReadPin()} & 读取按键或输入电平 \\
GPIO & \api{HAL_GPIO_WritePin()} & 控制 LED 或普通输出电平 \\
中断 & \api{HAL_NVIC_SetPriority()} & 设置中断抢占优先级和子优先级 \\
中断 & \api{HAL_NVIC_EnableIRQ()} & 使能 NVIC 中断通道 \\
中断 & \api{HAL_GPIO_EXTI_Callback()} & 外部中断回调，编写用户处理逻辑 \\
定时器 & \api{HAL_TIM_Base_Start_IT()} & 启动定时器周期中断 \\
定时器 & \api{HAL_TIM_PeriodElapsedCallback()} & 定时器更新中断回调 \\
UART & \api{HAL_UART_Transmit()} & 串口阻塞发送数据 \\
I2C & \api{HAL_I2C_Mem_Read()} & 从 I2C 器件寄存器读取数据 \\
SPI & \api{HAL_SPI_TransmitReceive()} & SPI 全双工收发数据 \\
DMA & \api{HAL_DMA_Start_IT()} & 启动 DMA 传输并开启中断 \\
ADC & \api{HAL_ADC_Start()} & 启动 ADC 转换 \\
ADC & \api{HAL_ADC_PollForConversion()} & 等待转换完成 \\
ADC & \api{HAL_ADC_GetValue()} & 读取 ADC 数字量 \\
ADC & \api{HAL_ADC_Start_IT()} & 中断方式启动 ADC \\
ADC & \api{HAL_ADC_Start_DMA()} & DMA 方式采样并搬运数据 \\
RTC & \api{HAL_RTC_SetTime()} & 设置时间 \\
RTC & \api{HAL_RTC_GetTime()} & 读取时间 \\
RTC & \api{HAL_RTC_SetAlarm_IT()} & 设置闹钟中断 \\
IWDG & \api{HAL_IWDG_Init()} & 初始化独立看门狗 \\
IWDG & \api{HAL_IWDG_Refresh()} & 喂狗，防止复位 \\
RCC & \api{HAL_RCC_OscConfig()} & 配置振荡源和 PLL \\
RCC & \api{HAL_RCC_ClockConfig()} & 配置系统时钟和总线分频 \\
SysTick & \api{HAL_GetTick()} & 获取毫秒节拍 \\
SysTick & \api{HAL_Delay()} & 毫秒延时 \\
UART & \api{HAL_UART_Receive_IT()} & 串口中断接收 \\
UART & \api{HAL_UART_RxCpltCallback()} & 串口接收完成回调 \\
GPIO & \api{HAL_GPIO_TogglePin()} & 翻转 LED 状态 \\
PWM & \api{HAL_TIM_PWM_Start()} & 启动 PWM 输出 \\
PWM & \api{__HAL_TIM_SET_COMPARE()} & 修改 CCR，改变占空比 \\
PWM & \api{HAL_TIM_PWM_Stop()} & 停止 PWM 输出 \\
\bottomrule
\end{longtable}
\end{center}

\section{考试速记}
\begin{summarybox}{一句话记忆}
\begin{itemize}
  \item ADC：模拟电压变数字量，核心公式是 $V=D/(2^N-1)\times V_{ref}$。
  \item RTC：掉电也能计时，优先记 LSE 32.768 kHz。
  \item 看门狗：程序要定期喂狗，超时复位；IWDG 用 LSI，WWDG 有窗口限制。
  \item 时钟树：HSI/HSE 进 PLL，倍频后得到 SYSCLK，再分到 AHB/APB。
  \item SysTick：内核 24 位定时器，常产生 1 ms 节拍。
  \item ADC 串口打印：启动、查询、读取、换算、格式化、发送。
  \item 按键 LED：按键读输入，LED 写输出，先判断高低电平有效。
  \item PWM：频率看 PSC 和 ARR，占空比看 CCR。
\end{itemize}
\end{summarybox}

\begin{summarybox}{简答题}
\textbf{题目：}利用通用定时器产生 PWM 信号驱动电机，用以调节智能小车的速度。要求信号的输出频率为 $10\text{ kHz}$，初始占空比为 $30\%$。若定时器输入时钟为 $72\text{ MHz}$，请根据题意设置预分频器 PSC、自动重装载寄存器 ARR 的值及比较寄存器 CCR 的值，并写出计算过程。\\[0.6em]
\textbf{解：}定时器 PWM 输出频率公式为：
\[
  f_{PWM}=\frac{f_{TIM}}{(PSC+1)(ARR+1)}
\]
已知 $f_{TIM}=72\text{ MHz}$，$f_{PWM}=10\text{ kHz}$，所以：
\[
  (PSC+1)(ARR+1)=\frac{72\times 10^6}{10\times 10^3}=7200
\]
为了计算方便，可令 $PSC=71$，则定时器计数频率为：
\[
  f_{CNT}=\frac{72\text{ MHz}}{PSC+1}=\frac{72\text{ MHz}}{72}=1\text{ MHz}
\]
因此：
\[
  ARR+1=\frac{1\text{ MHz}}{10\text{ kHz}}=100
\]
所以 $ARR=99$。初始占空比为 $30\%$，比较寄存器满足：
\[
  Duty=\frac{CCR}{ARR+1}\times 100\%
\]
因此：
\[
  CCR=(ARR+1)\times 30\%=100\times 0.3=30
\]
\textbf{答案：}可设置 $PSC=71$，$ARR=99$，$CCR=30$。此时 PWM 频率为 $10\text{ kHz}$，初始占空比为 $30\%$。\\[0.6em]
\textbf{说明：}本题答案不唯一，只要满足 $(PSC+1)(ARR+1)=7200$ 即可。
\end{summarybox}

\end{document}
